% ch-a4-06.tex —— 第 4 篇（Emerton–Gee–Hellmann）§5.3 续片，底本 PDF 240–251 = 印刷页 227–238

\ysrc{227}

……的定义，即映射 $\XX_P\to\XX_G$ 与 $\XX_P\to\XX_M$ 延拓为态射
\[
\begin{aligned}
\wt\beta_P&:\wt\XX_P^{\mathrm{naive}}\lra\XX_G,\\
\wt\alpha_P&:\wt\XX_P^{\mathrm{naive}}\lra\XX_M.
\end{aligned}
\]

此外还存在遗忘映射 $\wt\XX_P^{\mathrm{naive}}\to\ol\XX_P^{\mathrm{naive}}$：它在 $\XX_P$ 上诱导恒等映射，其定义是只记住 $V=V^{\check\lambda}$（$\check\lambda\in X^*(T)^{+}_{G}\cap X^*(T)_{G,P}$）所对应的映射 $\kappa_V$。此时 $M$-表示 $(V^{\check\lambda})^{U_P}$ 通过 $S_M$ 分解。

\begin{rem}{5.3.5}\label{rem:a4-5-3-5}
叠 $\ol\XX_P^{\mathrm{naive}}$ 与 $\wt\XX_P^{\mathrm{naive}}$ 在 $P=B$ 时重合，其他情形下则互不相同。在我们的设定 $G=\mathrm{GL}_d$ 中、且当 $P=B$ 时，只需记住 $\mathrm{GL}_d$ 的各基本权所对应的映射 $\kappa_{\check\lambda}$。也就是说，$\XX_B$ 的一个 $\mathrm{Sp}(A)$-值点可由元组 $(\mathcal E,\LL_1,\dots,\LL_d,\nu_1,\dots,\nu_d)$ 描述，其中 $\mathcal E$ 是秩 $d$ 的等变向量丛（equivariant vector bundle），$\LL_i$ 是等变线丛，$\nu_i$ 是态射
\[
\nu_i:\LL_i\lra\textstyle\bigwedge^{i}\mathcal E,
\]
它们在（相对于 $\mathrm{Sp}(A)$ 的）纤维上为单射、余核 $A$-平坦，且诸 $\nu_i$ 满足某些 Plücker 关系（我们不逐一写出）。
\end{rem}

结果表明，向量丛上的等变性条件带来一个新现象，使上述定义过于天真（「naive」之名即由此而来）：叠 $\XX_P$ 在 $\ol\XX_P^{\mathrm{naive}}$（相应地 $\wt\XX_P^{\mathrm{naive}}$）中并不稠密。粗略地说，这一现象的根源是 Hodge–Tate–Sen 权的存在。下面的例子把该现象展现得更具体：

\begin{exam}{5.3.6}\label{exam:a4-5-3-6}
设 $d=2$、$D\in\XX_2(L)$。则 $D$ 上的 $B$-结构由子对象 $\LL=\mathcal R_{F,L}(\delta_1)\subset D$ 给出，使得商 $D/\mathcal R_{F,L}(\delta_1)$ 是秩 $1$ 的 $(\varphi,\Gamma)$-模，记作 $\mathcal R_{F,L}(\delta_2)$。于是 $\mathrm{wt}_\tau(D)=\{\mathrm{wt}_\tau(\delta_1),\mathrm{wt}_\tau(\delta_2)\}$。由于 $D$ 与子对象 $\LL$ 的带标号 Hodge–Tate 权都在 $\ol\XX_B^{\mathrm{naive}}$ 上连续变化，故在 $\XX_B$ 的闭包上，子对象 $\LL$ 的权必是 $D$ 的诸权之一。另一方面，按 $\XX_B$ 的定义，对一切 $m\ge0$，对象 $(D,t^m\LL)$ 都在 $\XX_B$ 中。
\end{exam}

为了定义 $\XX_P$ 的良好紧化，我们希望以合理的方式取它在天真紧化中的闭包。当 $P=B$ 时，这正是下述定理的内容；该定理将在 [HHS] 中证明。

\begin{thm}{5.3.7}\label{thm:a4-5-3-7}
\begin{enumerate}[label=(\roman*),leftmargin=2.2em,itemsep=0.1em]
\item 叠 $\XX_B$ 与 $\ol\XX_B^{\mathrm{naive}}=\wt\XX_B^{\mathrm{naive}}$ 都是刚性解析 Artin 叠。
\item 存在良定义的 $\XX_B$ 在 $\ol\XX_B^{\mathrm{naive}}$ 中的「概形论」像 $\ol\XX_B$。
\end{enumerate}
\end{thm}

\begin{rem}{5.3.8}\label{rem:a4-5-3-8}
当然，我们预期对 $\XX_P$ 有类似的陈述，并由此以类似方式给出 $\XX_P$ 的定义。然而，从刚性解析空间到 $\XX_B$（相应地到 $\ol\XX_B^{\mathrm{naive}}$、$\ol\XX_B$）的光滑满射的构造，关键地用到 $\XX_T$ 是 Artin 叠这一事实。对更一般的抛物子群 $P$，只要假定 $\XX_M$ 是 Artin 叠，即能为 $M$ 证明猜想 5.1.19，我们就能对 $\XX_P$ 及其紧化 $\ol\XX_P$ 证明同样的结果。因此，对叠 $\XX_P$ 尚无完整的定义或刻画。下文我们仍用符号 $\XX_P$ 表示 $\XX_P$ 在 $\ol\XX_P^{\mathrm{naive}}$ 中的概形论像，它猜想性地良定义。类似地，人们期望\ysrc{228}通过取 $\XX_P$ 在 $\wt\XX_P^{\mathrm{naive}}$ 中的概形论像，得到紧化 $\wt\XX_P$。此时映射 $\alpha_P$、$\beta_P$、$\wt\alpha_P$、$\wt\beta_P$ 猜想性地分别限制到 $\ol\XX_P$、$\wt\XX_P$ 上，下文仍以同样的字母记之。
\end{rem}

\begin{proof}[证明梗概]
我们仅在二维情形勾勒「$\XX_B$ 是 Artin 叠」的证明。证明分两步：第一步，把 $\XX_B$ 限制到弱正则子集（weakly regular subset）$\XX_T^{\mathrm{wreg}}\subset\XX_T$ 上，得到 $\XX_B$ 的限制 $\XX_B^{\mathrm{wreg}}$。该弱正则子集保证所有 $\Ext^2$ 项消失，从而可以构造显式的卡。第二步，利用可以扭掉第二上同调 $H^2_{\varphi,\Gamma}$ 这一事实，并借助 Beauville–Laszlo 粘接引理，把一般情形化归到第一步。

\emph{第一步}：把满足 $H^2_{\varphi,\Gamma}(\mathcal R(\delta_1/\delta_2))=0$ 的特征 $(\delta_1,\delta_2)$ 构成的开子集记为 $\XX_T^{\mathrm{wreg}}\subset\XX_T$。若基域 $L$ 包含 $F$ 的全部伽罗瓦共轭，则 $\XX_T^{\mathrm{reg}}$ 可描述为一切满足下式的 $(\delta_1,\delta_2)$ 之集\tn{底本此处作 $\XX_T^{\mathrm{reg}}$，与上文 $\XX_T^{\mathrm{wreg}}$ 不一致，疑脱字母 w。}：
\[
\delta_1/\delta_2\notin\{\varepsilon z^\lambda\mid\lambda\in\textstyle\prod_{\tau:F\into L}\Z_+\}.
\]

设 $U=\mathrm{Sp}(A)$ 是 $\XX_T^{\mathrm{wreg}}$ 的仿射开子集。把 $(\delta_1,\delta_2)$ 记为万有特征在 $U$ 上的限制，并对 $\delta_i$ 沿从刚性空间到 $U$ 的映射的拉回沿用同一记号。由于 $H^2_{\varphi,\Gamma}(\mathcal R_{A,F}(\delta_1/\delta_2))$ 在 $U$ 上逐点消失，存在从集中在次数 $[0,1]$ 的完美复形到注 5.1.18 (i) 中计算 $H^\bullet_{\varphi,\Gamma}(\mathcal R_{A,F}(\delta_1/\delta_2))$ 的 Herr 复形 $C^\bullet_A$ 的拟同构
\[
[\mathcal E^0\xrightarrow{d}\mathcal E^1]\xrightarrow{f^\bullet}(C^\bullet_A,\partial^\bullet).
\]
而且该拟同构在任意拉回之后仍为拟同构。映射 $f^1:\mathcal E^1\to\ker\partial^1\subset C^1_A$ 给出与 $\mathcal E^1$ 相伴的几何向量丛 $Y\to U$ 上的万有扩张
\[
0\lra\mathcal R_{Y,F}(\delta_1)\lra D_Y\lra\mathcal R_{Y,F}(\delta_2)\lra 0.
\]
该扩张诱导态射
\[
g:Y\lra\XX_B|_U.
\]

我们断言 $g$ 是光滑满射。事实上，$g$ 按构造为满，又因 $\XX_B$ 的对角可表而相对可表。剩下只需验证 $g$ 满足无穷小提升准则。设 $B\onto\bar B$ 是仿射 $A$-代数的满射，考虑交换图
\[
\begin{tikzcd}[row sep=large,column sep=large]
\mathrm{Sp}(\bar B)\arrow[r,"\bar h"]\arrow[d] & Y\arrow[d,"g"]\\
\mathrm{Sp}(B)\arrow[r,"h"] & \XX_B.
\end{tikzcd}
\]
需找到提升 $\mathrm{Sp}(B)\to Y$ 使图交换。态射 $\bar h$ 定义了元素 $s_{\bar B}\in\Gamma(\mathrm{Sp}(\bar B),\mathcal E^1_{\bar B})$（其中 $\mathcal E^1_{\bar B}$ 记 $\mathcal E^1$ 到 $\mathrm{Sp}(\bar B)$ 的拉回，其余对象类似），其像 $f^1(s_{\bar B})$ 定义了扩张
\[
0\lra\mathcal R_{\bar B,F}(\delta_1)\lra D_{\bar B}\lra\mathcal R_{\bar B,F}(\delta_2)\lra 0.
\]
态射 $h$ 把扩张 $D_{\bar B}$ 提升为 $\mathcal R_{B,F}(\delta_2)$ 被 $\mathcal R_{B,F}(\delta_1)$ 的扩张 $D_B$；并且（在选定适当基之后）可把 $D_B$ 看作上闭链\ysrc{229}$c^1_B\in C^1_B$，它提升 $f^1(s_{\bar B})\in C^1_{\bar B}$。为找到提升 $\mathrm{Sp}(B)\to Y$，需找到元素 $s_B\in\Gamma(\mathrm{Sp}(B),\mathcal E^1_B)$，使由 $f^1(s_B)$ 定义的扩张与 $D_B$ 同构，且该同构提升 $D_{\bar B}$ 与由 $f^1(s_{\bar B})$ 定义的扩张之间的标准等同。换言之，需找到 $s_B$ 与在 $C^0_{\bar B}$ 中化为 $0$ 的 $c^0_B\in C^0_B$，使得
\[
c^1_B-f^1(s_B)=\partial^0(c^0_B).
\]
由于 $\mathcal E^1_B\to H^1(\mathcal E^\bullet_B)=H^1_{\varphi,\Gamma}(\mathcal R_{B,F}(\delta_1/\delta_2))$ 为满，存在 $s'_B\in\mathcal E^1_B$，它提升 $s_B$ 并代表扩张 $D_B$ 所定义的上同调类 $[c^1_B]$。于是存在 $c'^0_B\in C^0_B$ 使 $c^1_B-f^1(s'_B)=\partial^0(c'^0_B)$。我们需要修改 $s'_B$ 与 $c'^0_B$，因为先验地 $c'^0_B$ 未必在 $\bar B$ 上化为 $0$。然而由构造知，其归约在 $\partial^0$ 下的像为 $0$，即
\[
c'^0_{\bar B}\in\ker(\partial^0:C^0_{\bar B}\to C^1_{\bar B})=\ker(\mathcal E^0_{\bar B}\to\mathcal E^1_{\bar B}).
\]
特别地，可把 $c'^0_{\bar B}$ 看作 $\mathcal E^0_{\bar B}$ 的元素 $e^0_{\bar B}$。取 $c'^0_{\bar B}$ 的提升 $e^0_B\in\mathcal E^0_B$，并令
\[
\begin{aligned}
s_B&=s'_B+d(e^0_B),\\
c^0_B&=c'^0_B-f^0(e^0_B).
\end{aligned}
\]
于是 $c^0_{\bar B}=c'^0_{\bar B}-f^0(e^0_{\bar B})=0$，且
\[
c^1_B-f^1(s_B)=\partial^0(c'^0_B)-f^1(d(e^0_B))=\partial^0(c^0_B),
\]
正如所愿。

\emph{第二步}：设 $U\subset\XX_T$ 是某非弱正则点 $(\delta_1,\delta_2)$ 的开邻域，且存在从仿射空间 $\mathrm{Sp}(A)$ 出发的光滑满射。则存在 $N\gg0$ 使 $H^2_{\varphi,\Gamma}(\mathcal R_{A,F}(z^N\delta_1/\delta_2))=0$，参看注 5.1.18 (iii)。考虑 $(\delta'_1,\delta'_2)=(\delta_1,z^N\delta_2)\in\XX_T^{\mathrm{wreg}}$ 的邻域 $U'=\{(\eta_1,z^N\eta_2)\mid(\eta_1,\eta_2)\in U\}$，以及由 $(\eta_1,\eta_2)\mapsto(\eta_1,z^N\eta_2)$ 给出的态射 $a:U\to U'$。该态射嵌在下图中
\[
\begin{tikzcd}[row sep=large,column sep=large]
 & V'\arrow[d] & \\
 \XX_B|_U\arrow[r,"b"]\arrow[d] & \XX_B|_{U'}\arrow[d] & \\
 U\arrow[r,"a"] & U' &
\end{tikzcd}
\]
其中 $V'\to\XX_B|_{U'}$ 是从刚性解析空间出发的光滑满射（由第一步知其存在），而 $b$ 是拉回
\begin{equation}
\begin{tikzcd}[column sep=large,row sep=small]
0\arrow[r] & \mathcal R(\eta_1)\arrow[r]\arrow[d,equal] & D\arrow[r] & \mathcal R(\eta_2)\arrow[r] & 0\\
0\arrow[r] & \mathcal R(\eta_1)\arrow[r] & D'\arrow[r]\arrow[u,hook,"\iota"] & z^N\mathcal R(\eta_2)\arrow[r]\arrow[u,hook] & 0
\end{tikzcd}\tag{5.3.9}
\end{equation}
所定义的态射。

我们定义 $V\to\XX_B|_U$ 为纤维积
\[
V=V'\times_{\XX_B|_{U'}}\XX_B|_U\lra\XX_B|_U
\]
对第一个因子的投影。由于 $V'\to\XX_B|_{U'}$ 是光滑满射，剩下只需说明 $b$ 相对可表。这正是下文引理 5.3.10 的内容。至此证明了 $\XX_B$ 是 Artin 叠。

\ysrc{230}

关于如何把该策略用于证明 $\XX_B$ 与 $\ol\XX_B^{\mathrm{naive}}$ 的相应断言，参看下文注 5.3.11。
\end{proof}

\begin{lem}{5.3.10}\label{lem:a4-5-3-10}
把扩张 $D$ 送到由 (5.3.9) 定义的拉回 $D'$ 的态射 $b:\XX_B\to\XX_B$ 相对可表\tn{底本此处两端同为 $\XX_B$；据 p.229 的图示，应为 $\XX_B|_U\to\XX_B|_{U'}$，此处照录原文。}。
\end{lem}

\begin{proof}[证明梗概]
设 $V'=\mathrm{Sp}(A)\to\XX_B$ 是定义扩张 $D'$ 的态射。对 (5.3.9) 中任意的 $D$，态射 $\iota$ 诱导同构 $D'[1/t]\isom D[1/t]$。把标量扩张 $D\otimes_{\mathcal R_F}F_\infty[[t]]$ 记为 $\hat D$，$\hat D'$ 类似。由 Beauville–Laszlo 粘接引理，$(\varphi,\Gamma)$-模 $D$ 可由 $D'$ 以及自由 $(A\otimes_{\Q_p}F_\infty)((t))$-模 $\hat D'[1/t]$ 中的 $\Gamma$-稳定 $(A\otimes_{\Q_p}F_\infty)[[t]]$-格重建。可假定 $D'$ 自由，并选取 $D'$ 的一组基使第一个基向量张成子对象 $\mathcal R_{A,F}(\eta_1)$。用这组基，可把 $\hat D$ 看作仿射 Grassmann 流形（affine Grassmannian）
\[
\mathrm{Grass}_\infty=L(\mathrm{Res}_{F_\infty/\Q_p}\mathrm{GL}_2)/L^{+}(\mathrm{Res}_{F_\infty/\Q_p}\mathrm{GL}_2)
\]
的一个 $\Gamma$-稳定 $A$-值点。此时格 $\hat D'$ 与基点 $e_0\in\mathrm{Grass}_\infty$ 等同。此外，「$D$ 是 (5.3.9) 那样的扩张」这一事实等价于：格 $\hat D$ 落在闭子空间
\[
t^\nu L^{<0}U_\infty t^{-\nu}e_0\subset\mathrm{Grass}_\infty
\]
中，其中 $\nu=(1,-N)$，而 $U_\infty\subset\mathrm{Res}_{F_\infty/\Q_p}\mathrm{GL}_2$ 是幺幂上三角矩阵的 Weil 限制。反过来，$t^\nu L^{<0}U_\infty t^{-\nu}e_0$ 中的 $\Gamma$-稳定格恰好描述 $V'\to\XX_B$ 沿 $b$ 的拉回。

该证明有两处技术困难：
\begin{enumerate}[label=(\alph*),leftmargin=2.2em,itemsep=0.1em]
\item 因 $F_\infty$ 在 $\Q_p$ 上不有限，Weil 限制 $\mathrm{Res}_{F_\infty/\Q_p}\mathrm{GL}_2$ 不能由刚性解析空间表示。
\item 子空间 $t^\nu L^{<0}U_\infty t^{-\nu}e_0$ 是 ind-概形（或 ind-刚性解析空间），而非概形（或刚性解析空间）。
\end{enumerate}
稍微细致些，可通过下降到一个有限子扩张 $F_n=F(\mu_{p^n})\subset F_\infty$ 来克服 (a)。为克服 (b)，需说明 $\Gamma$-不变的格实际上落在 $t^\nu L^{<0}U_\infty t^{-\nu}e_0$ 的某个闭刚性解析子空间内。当 $F=\Q_p$ 时，该论证的主要成分如下（对一般的 $F$，同样的论证稍加变通即可）：在 $L$-值点上，或者 $\hat D$ 由 $\hat D'$ 与图 (5.3.9) 唯一确定，或者 $D'$ 在扭一下之后是 de Rham 的。在后一情形，$\hat D'$ 与 $\hat D$ 都落在仿射 Grassmann 流形中关于第三个基点 $\hat D''$ 的某些仿射 Schubert 胞腔内（在 de Rham 情形，可取格 $\hat D''$ 使其定义 Hodge–Tate 权为 $(0,0)$ 的 $(\varphi,\Gamma)$-模，而包含 $\hat D'$、$\hat D$ 的 Schubert 胞腔分别依赖于它们各自的 Hodge–Tate 权）。
\end{proof}

\begin{rem}{5.3.11}\label{rem:a4-5-3-11}
为证明 $\ol\XX_B^{\mathrm{naive}}$ 与 $\ol\XX_B$ 是刚性解析 Artin 叠，我们沿用上述梗概中的想法。为构造 $\ol\XX_B$ 的卡，须考虑 $t^\nu L^{<0}U_\infty t^{-\nu}e_0$ 在仿射 Grassmann 流形中的闭包；为构造 $\ol\XX_B^{\mathrm{naive}}$ 的卡，须考虑若干 $t^\nu L^{<0}U_\infty t^{-\nu}e_0$ 的并。
\end{rem}

\ysrc{231}

与代数曲线上向量丛的经典情形的一个重要差别是，映射
\[
\bar\beta_B:\ol\XX_B\lra\XX_G
\]
既不 proper 也不满（更一般的抛物子群也有类似现象）。此映射的像中的 $(\varphi,\Gamma)$-模都是三角化的（trianguline）：按 $\ol\XX_B^{\mathrm{naive}}$ 的定义，$\ol\XX_B^{\mathrm{naive}}$ 上（从而 $\ol\XX_B$ 上）的 $(\varphi,\Gamma)$-模在反转 $t$ 之后都容许一条完全的 $\varphi$ 与 $\Gamma$ 稳定旗；关于该旗在 Plücker 嵌入下的像，参看注 5.3.5。这说明 $\bar\beta_B$ 不满，因为并非每个伽罗瓦表示（更一般地，并非每个 $(\varphi,\Gamma)$-模）都是三角化的（等价地说：在反转 $t$ 之后都是三角化的）。看出这点最简便的办法是比较叠 $\XX_d$ 与 $\XX_B$ 的维数（或比较某些相关的伽罗瓦表示刚性解析空间的维数）。

态射 $\beta_B$ 也不 proper，下文注 5.3.15 将予以说明。不 proper 与不满是两个（乍看之下令人意外的）与代数曲线上向量丛叠之经典情形的差别。然而下述定理表明，$\bar\beta_B$ 仍相当接近 proper。

\begin{thm}{5.3.12}\label{thm:a4-5-3-12}
态射
\[
\bar\beta_B:\ol\XX_B\lra\XX_G\times\XX_T
\]
可由刚性解析空间表示，且为 proper。
\end{thm}

此结果利用 Liu [Liu15] 与 Kedlaya–Pottharst–Xiao [KPX14] 关于族中三角化的结果证得。当然我们预期紧化 $\ol\XX_P$ 与 $\wt\XX_P$ 也有类似结果，如下述猜想所述。我们还可以猜想性地把叠 $\XX_B$（更确切地说，它在 $\XX_d\times\XX_T$ 中的像）与三角化簇（trianguline variety）、从而与特征簇及 $p$-进自守形式联系起来。猜想 9.6.9 与 9.6.31 将使与 $p$-进自守形式空间的关系更为精确。

\begin{conj}{5.3.13}\label{conj:a4-5-3-13}
\begin{enumerate}[label=(\roman*),leftmargin=2.2em,itemsep=0.1em]
\item 态射
\begin{equation}
\begin{aligned}
\ol\XX_P&\lra\XX_G\times\XX_{S_M}\quad\text{与}\\
\wt\XX_P&\lra\XX_G\times\XX_M
\end{aligned}\tag{5.3.14}
\end{equation}
可由刚性解析空间表示，且为 proper。
\item 态射 (5.3.14) 的概形论像等维，维数为 $\dim\mathrm{Res}_{F/\Q_p}P$。
\item 设 $P=B$，并设 $R_{\bar\rho}$ 是固定剩余表示 $\bar\rho:\Gal_F\to\mathrm{GL}_d(k)$ 的加框架形变环（framed deformation ring）。则 (5.3.14) 的概形论像沿光滑映射
\[
(\Spf R_{\bar\rho})^{\mathrm{rig}}\times T^d\lra\XX_G\times\XX_T
\]
的拉回，与 [BHS17b] 的三角化簇重合。
\end{enumerate}
\end{conj}

这里「概形论像」自然应理解为刚性解析 Artin 叠背景下通常的概形论像概念的类比。

\ysrc{232}

\begin{rem}{5.3.15}\label{rem:a4-5-3-15}
尽管定理 5.3.12 或可解释为 $\bar\beta_B$ 非常接近 proper，但有一个方面它与 proper 映射截然不同：人们预期 $\XX_B$ 在 $\XX_G$ 中的像是 Zariski 稠密的（而 proper 映射的像是 Zariski 闭的，因为结构层的顺像是凝聚层）。这个像应看作 Gouvea–Mazur [GM98] 的「无限蕨」（infinite fern）及其由 Chenevier [Che13] 与 Nakamura [Nak14] 所作的推广的类比。无限蕨论证直接推出：像 $\bar\beta_B(\ol\XX_B)$ 在所有含内点晶体点的 $\XX_G$ 不可约分支之并中稠密。我们预期 $\XX_G$ 的所有分支都在其内部含有晶体点（作为关于 $\XX_G$ 连通分支的猜想性描述的推论，这实际上会是猜想 5.1.36 的推论）。

鉴于与特征簇的联系，不 proper 还可与下述事实关联：特征簇的构造涉及权空间上无限次的 Fredholm 超曲面。由于这里的不 proper 是在权空间上而非在伽罗瓦表示空间上不 proper，该蕴涵并非直接，但这两个现象是相关的。
\end{rem}

\begin{rem}{5.3.16}\label{rem:a4-5-3-16}
从纯形变论的观点看，等变 $P$-丛之叠 $\XX_P$ 的预期维数是 $\dim\mathrm{Res}_{F/\Q_p}P$。事实上我们预期 $\XX_P$ 等维且维数为此，并且是局部完全交。$p$-进情形似乎与 $L$-参数空间的情形颇为不同 [Hel23, Remark 2.2], [Zhu25a, Remark 2.3.8]：那里抛物子群的 $L$-参数叠可能不等维（而约化群的 $L$-参数叠总是等维的完全交）。因此在 $L$-参数叠的情形，由于这一现象，人们被迫在导出叠中工作。在 $p$-进情形，我们预期叠 $\XX_P$（以及类似地其紧化 $\ol\XX_P$ 与 $\wt\XX_P$）没有额外的非平凡导出结构。事实上，对固定的 $d$，可以证明当 $[F:\Q_p]$（与 $d$ 相比）很大时 $\XX_B$ 等维。目前我们还没有对一切 $F$ 的论证，但预期等维性与次数 $[F:\Q_p]$ 无关，似乎是合理的。
\end{rem}

受猜想 5.3.13 启发，把 $\XX_B$ 在 proper 映射 $\bar\beta_B$ 下的概形论像记为 $\XX_{G,\mathrm{tri}}\subset\XX_G\times\XX_T$。下述定理给出态射 $\XX_{G,\mathrm{tri}}\to\XX_G$ 的原像的部分描述，这些原像在第 9.6 节构造特征簇中的伴点（companion point，即在给定相伴伽罗瓦表示下构造有限斜率的超收敛 $p$-进自守形式）时将起重要作用。描述原像之前，需要一些准备。

设 $D\in\XX_G(L)$ 是晶体 $(\varphi,\Gamma)$-模，并设相伴 Weil–Deligne 表示 $\mathrm{WD}(D)$ 上 Frobenius 的特征值 $\varphi_1,\dots,\varphi_d$ 都在 $L$ 中，且对一切 $i\neq j$ 满足 $\varphi_i/\varphi_j\neq1,q$。此外，把 $D$ 的 Hodge–Tate 权记为
\[
\lambda=(\lambda_{\tau,1}\ge\cdots\ge\lambda_{\tau,d})_{\tau:F\into L},
\]
并设 $\lambda$ 正则，即对一切 $\tau$ 与 $i\neq j$ 有 $\lambda_{\tau,i}\neq\lambda_{\tau,j}$。固定 Frobenius 特征值 $\varphi_1,\dots,\varphi_d$ 的一个排序，等价地固定 $\mathrm{WD}(D)$ 上的一条完全 Frobenius 稳定旗 $\mathcal F_\bullet$。对每个嵌入 $\tau:F\into L$，把 Weyl 群元素 $w_{F,\tau}\in S_n$ 定义为 Hodge 滤过之 $\tau$-部分
\[
\mathrm{Fil}^\bullet_\tau\subseteq\mathrm{WD}(D)=\mathrm{WD}(D)_\tau\subseteq\textstyle\prod_{\tau':F\into L}\mathrm{WD}(D)_{\tau'}=\mathrm{WD}(D)\otimes_{\Q_p}L
\]
\ysrc{233}相对于 $\mathcal F_\bullet$ 的相对位置（注意由 $\lambda$ 的正则性假设，滤过 $\mathrm{Fil}^\bullet_\tau$ 确实是完全旗）\tn{底本此处作 $S_n$，而前文以 $d$ 记秩（$D$ 的 Hodge–Tate 权有 $d$ 个分量），疑为 $S_d$ 之误。}。

对旗 $\mathcal F_\bullet$（等价地，对 Frobenius 特征值 $\varphi_1,\dots,\varphi_d$ 的排序）配以非分歧特征 $\delta_{\mathcal F_\bullet}:T=(F^\times)^n\to L^\times$，它由
\begin{equation}
(x_1,\dots,x_n)\longmapsto\prod_{i=1}^{d}\mathrm{unr}_{\varphi_i}(x_i)\tag{5.3.17}
\end{equation}
给出\tn{底本在 $T=(F^\times)^n$ 中作 $n$，而右端求积指标到 $d$，二者应一致（$n=d$），此处照录原文。}。

此外，对一个元组 $(\mu_\tau)_\tau\in\Z^{\Hom(F,L)}$，有特征
\begin{equation}
z^\mu:x\longmapsto\prod_\tau\tau(x)^{\mu_\tau}.\tag{5.3.18}
\end{equation}
更一般地，对 $\underline\mu=(\mu_1,\dots,\mu_d)$（$\mu_i\in\Z^{\Hom(F,L)}$），记
\[
z^{\underline\mu}:(x_1,\dots,x_d)\longmapsto\prod_{i=1}^{d}z^{\mu_i}(x_i)=\prod_{i=1}^{d}\prod_\tau\tau(x_i)^{\mu_{i,\tau}}.
\]

$\XX_{G,\mathrm{tri}}$ 中位于 $D$ 上方的点的下述刻画来自 [BHS19, Theorem 4.2.3]。

\begin{thm}{5.3.19}\label{thm:a4-5-3-19}
设 $D\in\XX_G(L)$ 是 Hodge–Tate 权 $\lambda$ 正则的晶体 $(\varphi,\Gamma)$-模，且 $\mathrm{WD}(D)$ 上 Frobenius 的特征值 $\varphi_1,\dots,\varphi_d$ 都在 $L$ 中，并对 $i\neq j$ 满足 $\varphi_i/\varphi_j\neq1,q$。则
\[
\begin{aligned}
&\{\delta\in\XX_T\mid(D,\delta)\in\XX_{G,\mathrm{tri}}\}\\
&\quad=\bigcup_{\mathcal F_\bullet}\{\delta_{\mathcal F_\bullet}z^{w\lambda}\mid w=(w_\tau)_\tau\in W\ \text{使得对一切 }\tau\text{ 有 }w_{F,\tau}\preceq w_\tau w_0\},
\end{aligned}
\]
其中并集遍历 $\mathrm{WD}(D)$ 的一切 Frobenius 稳定旗 $\mathcal F$，而 $w_0\in S_d$ 是最长元素。
\end{thm}

\begin{rem}{5.3.20}\label{rem:a4-5-3-20}
\begin{enumerate}[label=(\roman*),leftmargin=2.2em,itemsep=0.1em]
\item 沿用 Hansen [Han17, Conjecture 6.2.3]，对 $\XX_{G,\mathrm{tri}}\to\XX_G$ 的纤维有一般的猜想性描述，无需假定 $(\varphi,\Gamma)$-模 $D$ 晶体（也不需要定理中的正则性假设）。
\item 对 crystabelline $(\varphi,\Gamma)$-模，该描述与晶体情形大致相同。在半稳定情形（更一般地，在 semistabelline 情形），只取那些在 Weil–Deligne 表示的单值作用下稳定的旗 $\mathcal F_\bullet$。
\item 由 semistabelline 情形推广到一切 de Rham $(\varphi,\Gamma)$-模 $D$，是合理的。此时 $\XX_B$ 在 $\XX_G\times\XX_T$ 中的像应换成 $\XX_P$ 在 $\XX_G\times\XX_{S_M}$ 中的像，其中 $P$ 是包含 $B$ 且 Levi 为 $M=\prod_{i=1}^{r}\mathrm{GL}_{d_i}$ 的抛物子群，使得 $D$ 对应的 Weil–Deligne 表示 $\mathrm{WD}(D)$ 是不可分解 Weil–Deligne 表示 $D_1,\dots,D_r$ 的直和，其维数分别为 $d_i$（参看 [BD24]）。
\end{enumerate}
\end{rem}

下述预期源于三角化簇与特征簇、以及特征簇的拼接版本 [BHS17b] 之间的联系（也源于猜想 6.2.4(2) 与猜想 9.6.18）。

\ysrc{234}

\begin{conj}{5.3.21}\label{conj:a4-5-3-21}
把标准（proper）投影记为 $\pi_B:\ol\XX_B\to\XX_G\times\XX_T$。
\begin{enumerate}[label=(\roman*),leftmargin=2.2em,itemsep=0.1em]
\item 有同构
\[
R\pi_{B,*}\OO_{\ol\XX_B}\cong\OO_{\XX_{G,\mathrm{tri}}},\qquad
R\pi_{B,*}\omega_{\ol\XX_B}\cong\omega_{\XX_{G,\mathrm{tri}}}.
\]
\item 顺像 $R\pi_{B,*}(\OO_{\ol\XX_B}([F:\Q_p]\rho'))$ 集中在次数 $0$，且为 Cohen–Macaulay 模（即其对偶 $R\pi_{B,*}(\omega_{\ol\XX_B}(-[F:\Q_p]\rho'))$ 也集中在次数 $0$）。
\end{enumerate}
\end{conj}

此处 $\rho'$ 是 $T=\GG_m^d$ 的代数特征 $(0,-1,\dots,-d+1)$，我们把它看作 $*/\GG_m^d$ 上的线丛，并沿标准映射
\[
\ol\XX_B\lra\XX_T=\XX_1^d=(T/\GG_m)^d\lra*/\GG_m^d
\]
把它拉回到 $\ol\XX_B$ 上；关于 $\XX_1$ 的辨认见第 7.1 节。

此外，追问三角化簇 $\XX_{G,\mathrm{tri}}$ 是否正规、是否 Cohen–Macaulay 是有意义的（注意若 $\ol\XX_B$ 是 Cohen–Macaulay，则由猜想的 (i) 可推出 Cohen–Macaulay 性）。我们指出，猜想 (ii) 中的层应看作 [BZCHN24] 中凝聚 Springer 层的变体（相应地，[Zhu25a, 4.4] 或 [Hel23, 2.3] 的变体），因此 (ii) 应看作 [BZCHN24, Conjecture 4.15] 的变体（相应地，[Zhu25a, Conjecture 4.4.2] 或 [Hel23, Conjecture 2.17] 的变体）。

当 $F=\Q_p$、$G=\mathrm{GL}_2$ 时，我们可以证明上述猜想，并算出 (ii) 中 Cohen–Macaulay 模的奇点与纤维维数（这些纤维维数在「局部解析表示与 $\XX_d$ 上凝聚层的猜想性对应」（猜想 6.2.4）中有某种解释）。

\begin{thm}{5.3.22}\label{thm:a4-5-3-22}
设 $F=\Q_p$、$d=2$，即 $G=\mathrm{GL}_2$。
\begin{enumerate}[label=(\roman*),leftmargin=2.2em,itemsep=0.1em]
\item 叠 $\ol\XX_B$ 光滑。
\item 叠 $\XX_{G,\mathrm{tri}}$ 是刚性解析 Artin 叠\footnote{由于我们不知道 $\XX_G=\XX_2$ 是否为刚性解析 Artin 叠，此处应理解为：存在刚性解析 Artin 叠 $\XX_{G,\mathrm{tri}}$，它作为闭子叠嵌入 $\XX_G\times\XX_T$。映射 $\ol\XX_B\to\XX_G\times\XX_T$ 经由 $\XX_{G,\mathrm{tri}}$ 分解，并且概形论稠密。}，且正规、Cohen–Macaulay。
\item 态射 $\pi_B:\ol\XX_B\to\XX_{G,\mathrm{tri}}$ 双有理，纤维几何连通，且标准映射
\[
\OO_{\XX_{G,\mathrm{tri}}}\lra R\pi_{B,*}\OO_{\ol\XX_B}
\]
为同构；由对偶性还可得同构
\[
R\pi_{B,*}\omega_{\ol\XX_B}\cong\omega_{\XX_{G,\mathrm{tri}}}.
\]
\item 令 $\mathcal F=R\pi_{B,*}(\OO_{\ol\XX_B}(\rho'))$、$\mathcal G=R\pi_{B,*}(\omega_{\ol\XX_B}(-\rho'))$。则 $\mathcal F$ 与 $\mathcal G$ 集中在次数 $0$，特别地都是 Cohen–Macaulay 模。
\item \ysrc{235}对 $x=(D,\delta_1,\delta_2)\in\XX_{G,\mathrm{tri}}$，纤维 $\mathcal F\otimes k(x)$、$\mathcal G\otimes k(x)$ 与 $\omega_{\XX_{G,\mathrm{tri}}}\otimes k(x)$ 都是一维的，除非 $D=\mathcal R(\delta)\oplus t^{-n}\mathcal R(\delta)$（$n\ge0$）且 $\delta_1=\delta$。此时
\[
\begin{aligned}
\dim\mathcal F\otimes k(x)&=2,\\
\dim\mathcal G\otimes k(x)&=\begin{cases}2,&n=0\\1,&n\ge1\end{cases},\\
\dim\omega_{\XX_{G,\mathrm{tri}}}\otimes k(x)&=\begin{cases}3,&n=0\\2,&n\ge1.\end{cases}
\end{aligned}
\]
\end{enumerate}
\end{thm}

\medskip\noindent\textbf{5.3.23. 局部模型.} 本节余下部分假定 $L$ 包含 $F$ 的全部伽罗瓦共轭（在 $\bar F$ 内）。我们讨论 $\mathcal R_F$ 上 $(\varphi,\Gamma)$-模之叠与 $F_\infty[[t]]$ 上（半线性）$\Gamma$-表示之叠的关系。我们不知道是否应期待后者是 Artin 叠（它们似乎可以相当病态），但它们在固定 Hodge–Tate 权处的完备化表现良好。在某些「好」点上，从 $(\varphi,\Gamma)$-模之叠到这些叠的态射形式光滑，因此我们得到 $\ol\XX_B$ 之类叠的相当显式的局部模型，推广了 [BHS19] 中三角化簇的局部模型。

按如下方式在 $\mathrm{Rig}_L$ 上定义群胚 $\XX_d^{\mathrm{dR},+}$。对仿射代数 $A$，令
\[
\XX_d^{\mathrm{dR},+}(\mathrm{Sp}(A))=\left\{\begin{array}{l}
\text{秩 }d\text{ 的有限投射 }A\hat\otimes F_\infty[[t]]\text{-模上的}\\
\text{连续、半线性 }\Gamma\text{-表示}
\end{array}\right\},
\]
其中（记号略作滥用）$A\hat\otimes F_\infty[[t]]=(A\otimes_{\Q_p}F_\infty)[[t]]$。

\begin{rem}{5.3.24}\label{rem:a4-5-3-24}
\begin{enumerate}[label=(\roman*),leftmargin=2.2em,itemsep=0.1em]
\item 类似地，把 $\mathrm{Sp}(A)$ 送到「秩 $d$ 的有限投射 $A\hat\otimes F_\infty((t))$-模附带半线性 $\Gamma$-作用，且（在 $\mathrm{Sp}(A)$ 上局部）容许一个 $A\hat\otimes F_\infty[[t]]$-格」构成的群胚，以此定义群胚 $\XX_d^{\mathrm{dR}}$。
\item 考虑这些叠的动机（也是其名称的由来）是 Fontaine 的定理 [Fon04]：有限自由 $F_\infty[[t]]$-模（相应地，有限维 $F_\infty((t))$-向量空间）上全体半线性 $\Gamma$-表示构成的范畴，等价于有限自由 $B^+_{\mathrm{dR}}$-模（相应地，有限维 $B_{\mathrm{dR}}$-向量空间）上全体半线性 $\Gal_F$-表示构成的范畴。我们预期类似的等价在刚性解析空间上的族中也成立；但由于尚无证明，我们只关注半线性 $\Gamma$-表示，它们似乎更易处理。
\end{enumerate}
\end{rem}

\medskip\noindent\textbf{期望 5.3.25.}\label{exp:a4-5-3-25}
\begin{enumerate}[label=(\roman*),leftmargin=2.2em,itemsep=0.1em]
\item 群胚 $\XX_d^{\mathrm{dR},+}$ 与 $\XX_d^{\mathrm{dR}}$ 都是叠。
\item $\XX_d^{\mathrm{dR},+}$ 的对角可表。
\item 态射 $\XX_d^{\mathrm{dR},+}\to\XX_d^{\mathrm{dR}}$ 可表。
\end{enumerate}

尽管我们不知道是否应期待 $\XX_d^{\mathrm{dR},+}$ 与 $\XX_d^{\mathrm{dR}}$ 是 Artin 叠，这两个叠在刚性解析点处都有万有形变环；更一般地，沿固定 Hodge–Tate–Sen 权的完备化是形式刚性解析 Artin 叠（此概念我们不在此正式引入）。为使该陈述有意义，我们指出：与 (5.1.35) 中权映射 $\omega_d$ 的定义类似，\ysrc{236}有权映射
\[
\omega_d^{\mathrm{dR}}:\XX_d^{\mathrm{dR},+}\lra\mathrm{WT}_{d,F}\cong\textstyle\prod_{F\into L}\mathbb A^d/S_d,
\]
它在 $A$-值点上把一个 $A\hat\otimes F_\infty[[t]]$-模 $D_A$ 送到 $\Gamma$-作用在 $D_A/tD_A$ 上的导数的特征多项式（回忆我们已取 $L$ 足够大，使 $[F:\Q_p]=|\Hom(F,L)|$）。

\begin{rem}{5.3.26}\label{rem:a4-5-3-26}
对象 $D_A\in\XX_d^{\mathrm{dR}}(\mathrm{Sp}(A))$ 的 Hodge–Tate–Sen 权只在模 $\Z$ 的意义下良定义。更精确地说，在同构 $A\otimes_{\Q_p}F=\prod_{F\into L}A$ 之下，可把 $D_A$ 的权定义为
\[
\textstyle\prod_{F\into L}\big((\mathbb A/\Z)^d/S_d\big)
\]
的元素。
\end{rem}

Fontaine 关于几乎 de Rham（almost de Rham）表示的理论给出叠 $\XX_d^{\mathrm{dR}}$ 与 $\XX_d^{\mathrm{dR},+}$ 在固定整 Hodge–Tate 权处完备化后的描述，且在任意固定 Hodge–Tate–Sen 权处的完备化也有类似结果。我们简要回顾 Fontaine 对 $B_{\mathrm{dR}}$-表示（等价地，$F_\infty((t))$-表示）的分类 [Fon04]。把有限维 $B_{\mathrm{dR}}$ 向量空间上连续半线性 $\Gal_F$-表示构成的范畴记为 $\mathrm{Mod}^{\Gal_F}_{B_{\mathrm{dR}}}$，并类似地以显然的方式定义 $\mathrm{Mod}^{\Gal_F}_{B^+_{\mathrm{dR}}}$、$\mathrm{Mod}^\Gamma_{F_\infty((t))}$ 与 $\mathrm{Mod}^\Gamma_{F_\infty[[t]]}$。对 Hodge–Tate–Sen 权 $a\in F$，把 Hodge–Tate–Sen 权模 $\Z$ 同余于 $a$ 的一切对象构成的全子范畴记为 $\mathrm{Mod}^{\Gal_F,a}_{B_{\mathrm{dR}}}$。则
\[
\mathrm{Mod}^{\Gal_F}_{B_{\mathrm{dR}}}=\bigoplus_{a\in F/\Z}\mathrm{Mod}^{\Gal_F,a}_{B_{\mathrm{dR}}}
\]
是 $\mathrm{Mod}^{\Gal_F}_{B_{\mathrm{dR}}}$ 的块分解，即不同直和项中对象之间的高阶 $\Ext$-群消失（这是 [Fon04, Thm. 3.19] 的推论）。

块 $\mathrm{Mod}^{\Gal_F,0}_{B_{\mathrm{dR}}}$ 称为几乎 de Rham 表示之块（法语作 presque de Rham）。用线性代数对象可把它描述如下（见 [BHS19, Prop 3.1.1]）：

\begin{thm}{5.3.27 (Fontaine)}\label{thm:a4-5-3-27}
存在范畴等价
\[
D_{\mathrm{pdR}}:\mathrm{Mod}^{\Gal_F,0}_{B^+_{\mathrm{dR}}}\lra\mathrm{FilRep}(\GG_a),
\]
从 $B^+_{\mathrm{dR}}$ 上几乎 de Rham 表示构成的范畴，到「有限维 $F$-向量空间附带加法群 $\GG_a$ 的代数表示，以及一条分离、穷尽、且在 $\GG_a$-作用下稳定的 $\Z$-滤过 $\mathrm{Fil}^\bullet$」构成的范畴。在该等价下，从 $\mathrm{Mod}^{\Gal_F,0}_{B^+_{\mathrm{dR}}}$ 过渡到 $\mathrm{Mod}^{\Gal_F,0}_{B_{\mathrm{dR}}}$ 对应于忘掉滤过。
\end{thm}

我们回忆：在特征 $0$ 的域上，$\GG_a$-表示不过是幺幂算子的作用。此外，$B^+_{\mathrm{dR}}$ 上（或 $B_{\mathrm{dR}}$ 上）的几乎 de Rham 表示 $V$ 是 de Rham 的，当且仅当幺幂算子 $N:D_{\mathrm{pdR}}(V)\to D_{\mathrm{pdR}}(V)$（由 $\GG_a$-作用定义）为零。

函子 $D_{\mathrm{pdR}}$ 有周期环风格的刻画 [BHS19, 3.1]，我们不在此回顾。在等价范畴 $\mathrm{Mod}^\Gamma_{F_\infty[[t]]}\cong\mathrm{Mod}^{\Gal_F}_{B^+_{\mathrm{dR}}}$ 上，几乎 de Rham 表示与函子 $D_{\mathrm{pdR}}$ 可刻画如下：给定带连续 $\Gamma$-作用的 $F_\infty[[t]]$-模 $\Lambda$，可以考虑\ysrc{237}（与 Hodge–Tate–Sen 权的定义类似）$\Gamma$-作用在 $1$ 处的导数，它给出一个导数
\[
\partial_\Lambda:\Lambda\lra\Lambda
\]
位于 $F_\infty[[t]]$ 上的导数 $t\frac{d}{dt}$ 之上。于是可定义
\[
D_{\mathrm{pdR},\infty}(\Lambda)=\bigcup_{N\ge1}\ker(\partial_\Lambda^N:\Lambda[1/t]\to\Lambda[1/t]),
\]
即 $\Lambda[1/t]$ 中 $\partial_\Lambda$ 幺幂的那个子 $F_\infty$-向量空间。特别地，$\partial_\Lambda$ 诱导 $D_{\mathrm{pdR},\infty}(\Lambda)$ 上的 $F_\infty$-线性幺幂自同态 $N_\infty$。表示 $\Lambda$ 几乎 de Rham，当且仅当标准映射
\[
D_{\mathrm{pdR},\infty}(\Lambda)\otimes_{F_\infty}F_\infty((t))\lra\Lambda[1/t]
\]
为同构。此时记 $\Lambda_0=D_{\mathrm{pdR},\infty}(\Lambda)\otimes_{F_\infty}F_\infty[[t]]$，并称 $\Lambda_0$ 为标准格（standard lattice）。于是格 $\Lambda\subset\Lambda_0[1/t]$ 通过子 $F_\infty$-向量空间给出 $D_{\mathrm{pdR},\infty}(\Lambda)=\Lambda_0/t\Lambda_0$ 上的一条分离、穷尽的 $\Z$-滤过 $\mathrm{Fil}^\bullet_\infty$，这些子空间在 $N_\infty$ 的作用下稳定，因为 $\Lambda$ 在 $\partial_\Lambda$ 下稳定。同样与 Hodge–Tate–Sen 权的定义类似，三元组 $(D_{\mathrm{pdR},\infty}(\Lambda),N_\infty,\mathrm{Fil}^\bullet_\infty)$ 带有 $\Gamma$-作用，使我们能把 $(D_{\mathrm{pdR},\infty}(\Lambda),N_\infty,\mathrm{Fil}^\bullet_\infty)$ 从 $F_\infty$ 下降\footnote{提醒读者：该下降并不完全形式化，需要一些工作。回忆在讨论 Hodge–Tate 权时我们已指出，$(\mathrm{D}_{\mathrm{Sen}}(D_A),\Theta)$ 不能下降为 $F\otimes_{\Q_p}A$ 上的对象。然而在固定整 Hodge–Tate 权处完备化之后可以做到。此处的问题类似。}为 $F$ 上所要的三元组 $(D_{\mathrm{pdR}}(\Lambda),N,\mathrm{Fil}^\bullet)$。

Fontaine 定理中的等价是下述关于 $B^+_{\mathrm{dR}}$-表示之叠的完备化描述的核心。注意对 $L$-代数 $A$ 有
\[
F_\infty[[t]]\hat\otimes_{\Q_p}A=\prod_{\tau:F\into L}F_\infty[[t]]\hat\otimes_F A,
\]
因此可对每个嵌入 $\tau$ 与每个带标号 Hodge–Tate–Sen 权分别应用 Fontaine 的描述。本节余下部分记 $G=\mathrm{Res}_{F/\Q_p}\mathrm{GL}_d$、$\frakg=\mathrm{Lie}(G)$。此外固定一个整的带标号 Hodge–Tate 权
\[
\lambda\in\prod_{\tau:F\into L}(\Z^d)/S_d\subseteq\mathrm{WT}_{d,F}(L).
\]
用 $[\lambda]$ 记它模 $\Z$ 的类（它自然只记住该权是整的，并不区分不同的 $\lambda$）。我们描述群胚 $\XX_d^{\mathrm{dR},+}$ 与 $\XX_d^{\mathrm{dR}}$ 沿固定权 $\lambda$（相应地沿 Hodge–Tate–Sen 权模 $\Z$ 同余于 $\lambda$ 的轨迹）的形式完备化 $(\XX_d^{\mathrm{dR},+})^{\wedge}_{[\lambda]}$ 与 $(\XX_d^{\mathrm{dR}})^{\wedge}_{[\lambda]}$。该结果与 [BHS19, Cor. 3.1.6, Cor. 3.1.9] 非常相似。令
\[
\wt{\frakg}_{P_\lambda}=\{(A,\mathcal F^\bullet)\in\frakg\times G/P_\lambda\mid A\mathcal F^\bullet\subset\mathcal F^\bullet\}\subseteq\frakg\times G/P_\lambda
\]
为 Grothendieck 奇点同时消解的这一变体，它参数化由自同态 $A$ 与 $A$-稳定的 $P_\lambda$ 型旗 $\mathcal F^\bullet$ 构成的配对 $(A,\mathcal F^\bullet)$（记号参看 §5.2）。回忆 $\wt{\frakg}_{P_\lambda}$ 是 $\prod_\tau\mathfrak{gl}_d\times\mathrm{GL}_d/P_{\lambda_\tau}$ 的子簇，并且我们仍把 $G/P_\lambda$ 上的滤过看作一列在次数 $\lambda$ 处跳跃的 $\Z$-滤过所构成的 $[F:\Q_p]$-元组。

\ysrc{238}

\begin{prop}{5.3.28}\label{prop:a4-5-3-28}
函子 $D_{\mathrm{pdR},\infty}$ 诱导同构
\[
(\XX_d^{\mathrm{dR}})^{\wedge}_{[\lambda]}\isom\wt{\frakg}_0/G,\qquad
(\XX_d^{\mathrm{dR},+})^{\wedge}_{\lambda}\isom(\wt{\frakg}_{P_\lambda})^{\wedge}_0/G,
\]
其中右端在各幺幂自同态构成的闭子空间处完备化（即分别在幂零锥、及其部分 Springer 消解处完备化）。
\end{prop}

\begin{rem}{5.3.29}\label{rem:a4-5-3-29}
\begin{enumerate}[label=(\alph*),leftmargin=2.2em,itemsep=0.1em]
\item 该命题说明 $\XX_d^{\mathrm{dR},+}$ 在整 Hodge–Tate 权处的完备化是形式刚性解析 Artin 叠（对此类对象需作适当定义）。
\item 事实上，对任意（即不必整的）权 $\lambda$ 应用 Fontaine 的理论，可得到 $\XX_d^{\mathrm{dR},+}$（相应地 $\XX_d^{\mathrm{dR}}$）在 $\lambda$ 处完备化的类似描述。
\end{enumerate}
\end{rem}

在某些「好」点的邻域内，可把 $\XX_d^{\mathrm{dR},+}$ 看作叠 $\XX_d$ 的局部模型。下述定理是这方面的一个样本结果，可与 [HMS22, Theorem 3.1.1] 比较。下文还将讨论 $B$-结构（相应地 $P$-结构）及其 Drinfeld 紧化的相关情形。

\begin{thm}{5.3.30}\label{thm:a4-5-3-30}
设 $x\in\XX_d$ 是 Hodge–Tate 权正则的晶体点，其 Frobenius 特征值 $\varphi_1,\dots,\varphi_d$ 满足 $\varphi_i/\varphi_j\notin\{1,q\}$。则态射 $\XX_d\to\XX_d^{\mathrm{dR},+}$ 在 $x$ 处形式光滑。
\end{thm}

给定抛物子群 $P\subseteq\mathrm{GL}_d$，可考虑 $\XX_d^{\mathrm{dR}}$ 与 $\XX_d^{\mathrm{dR},+}$ 中对象上的 $P$-结构（即部分旗），从而定义叠 $\XX_P^{\mathrm{dR}}$ 与 $\XX_P^{\mathrm{dR},+}$。与上述关于 Drinfeld 紧化的讨论平行，我们希望定义 $\XX_P^{\mathrm{dR},+}$ 的「紧化」$\ol\XX_P^{\mathrm{dR},+}$ 与 $\wt\XX_P^{\mathrm{dR},+}$，使它们嵌入下图
\[
\begin{tikzcd}[row sep=large,column sep=large]
 & \wt\XX_P^{\mathrm{dR},+}\arrow[dl]\arrow[d]\arrow[dr] & \\
 \XX_d^{\mathrm{dR},+} & \ol\XX_P^{\mathrm{dR},+}\arrow[l]\arrow[dr] & \XX_M^{\mathrm{dR},+}\arrow[d]\\
 & & \XX_{S_M}^{\mathrm{dR},+}
\end{tikzcd}
\]
我们将在固定整权 $\lambda$ 处形式完备化之后定义这些紧化，并证明关于它们的有限性陈述。

如上，固定一个整的带标号 Hodge–Tate 权
\[
\lambda=((\lambda_{\tau,1},\dots,\lambda_{\tau,d})_{\tau:F\into L}.
\]
我们只限于最简单的情形：$\lambda$ 正则。由于 $\lambda$ 只在置换意义下良定义，取其支配代表元，即假定对一切 $\tau$ 有 $\lambda_{\tau,1}>\lambda_{\tau,2}>\cdots>\lambda_{\tau,d}$。特别地，抛物子群 $P_\lambda$ 等于 Borel 子群的 Weil 限制 $(\mathrm{Res}_{F/\Q_p}B)_L$。为简单起见记 $\wt{\frakg}=\wt{\frakg}_{P_\lambda}$\tn{底本此处左括号双重、右括号单一（$\lambda=((\lambda_{\tau,1},\dots,\lambda_{\tau,d})_{\tau:F\into L}.$），括号不配对，此处照录。}。

% 本片止于印刷页 238 页末（§5.3），续片自 PDF 252 页起
